\documentclass[12pt,oneside]{book}

% =====================================================
% METADATOS
% =====================================================
\newcommand{\universidad}{Universidad de Buenos Aires}
\newcommand{\facultad}{Facultad de Derecho}
\newcommand{\materia}{Derecho Procesal Civil}
\newcommand{\titulo}{Unidad 15 --- Prueba Pericial, Prueba Indiciaria y Reconocimiento Judicial}
\newcommand{\autor}{Isaac Edgar Camacho Ocampo}
\newcommand{\anio}{2020}

% =====================================================
% PAQUETES
% =====================================================
\usepackage[utf8]{inputenc}
\usepackage[T1]{fontenc}
\usepackage[spanish,es-nodecimaldot]{babel}

\usepackage[
    top=2.5cm,
    bottom=2.5cm,
    left=3cm,
    right=2.5cm,
    headheight=15pt
]{geometry}

\usepackage{amsmath}
\usepackage{amssymb}
\usepackage{mathtools}

\usepackage{graphicx}
\usepackage{float}
\usepackage{caption}
\usepackage{subcaption}

\usepackage{booktabs}
\usepackage{array}
\usepackage{longtable}
\usepackage{tabularx}

\usepackage{enumitem}

\usepackage{xcolor}
\usepackage[most]{tcolorbox}

\usepackage{fancyhdr}
\usepackage{ragged2e}

\usepackage[
    colorlinks=true,
    linkcolor=blue!50!black,
    citecolor=green!40!black,
    urlcolor=blue!60!black,
    pdftitle={\titulo},
    pdfauthor={\autor},
    pdfsubject={Apuntes universitarios},
    bookmarks=true
]{hyperref}

% =====================================================
% CONFIGURACIÓN
% =====================================================
\graphicspath{
    {./img/}
    {./graficos/}
    {./figuras/}
}

\setcounter{tocdepth}{2}

\setlength{\parindent}{0pt}
\setlength{\parskip}{0.5em}

\setlist[itemize]{
    topsep=0.3em,
    itemsep=0.2em
}

\setlist[enumerate]{
    topsep=0.3em,
    itemsep=0.2em
}

\clubpenalty=10000
\widowpenalty=10000

% =====================================================
% COMANDOS
% =====================================================
\newcommand{\abs}[1]{\left|#1\right|}
\newcommand{\norm}[1]{\left\lVert#1\right\rVert}
\newcommand{\vect}[1]{\mathbf{#1}}
\newcommand{\deriv}[2]{\frac{d#1}{d#2}}
\newcommand{\pderiv}[2]{\frac{\partial#1}{\partial#2}}

% =====================================================
% ENTORNOS / CAJAS
% =====================================================
\newtcolorbox{definition}[1]{
    enhanced,
    breakable,
    title={Definición: #1},
    fonttitle=\bfseries,
    colback=blue!3,
    colframe=blue!50!black,
    boxrule=0.8pt,
    arc=2mm
}

\newtcolorbox{important}{
    enhanced,
    breakable,
    title={Importante},
    fonttitle=\bfseries,
    colback=yellow!8,
    colframe=orange!70!black,
    boxrule=0.8pt,
    arc=2mm
}

\newtcolorbox{example}[1]{
    enhanced,
    breakable,
    title={Ejemplo: #1},
    fonttitle=\bfseries,
    colback=green!4,
    colframe=green!40!black,
    boxrule=0.8pt,
    arc=2mm
}

\newtcolorbox{formula}[1]{
    enhanced,
    breakable,
    title={Fórmula / Regla: #1},
    fonttitle=\bfseries,
    colback=purple!4,
    colframe=purple!50!black,
    boxrule=0.8pt,
    arc=2mm
}

\newtcolorbox{exercise}[1]{
    enhanced,
    breakable,
    title={Ejercicio: #1},
    fonttitle=\bfseries,
    colback=gray!5,
    colframe=gray!60!black,
    boxrule=0.8pt,
    arc=2mm
}

\newtcolorbox{note}{
    enhanced,
    breakable,
    title={Nota},
    fonttitle=\bfseries,
    colback=white,
    colframe=gray!50,
    boxrule=0.6pt,
    arc=2mm
}

% =====================================================
% CONFIGURACIÓN FINAL (encabezados y pies)
% =====================================================
\pagestyle{fancy}
\fancyhf{}

\fancyhead[L]{\nouppercase{\leftmark}}
\fancyhead[R]{\materia}

\fancyfoot[C]{\thepage}

\renewcommand{\headrulewidth}{0.4pt}
\renewcommand{\footrulewidth}{0pt}

\fancypagestyle{plain}{
    \fancyhf{}
    \fancyfoot[C]{\thepage}
    \renewcommand{\headrulewidth}{0pt}
}

% =====================================================
% DOCUMENTO
% =====================================================
\begin{document}

% -----------------------------------------------------
% PORTADA
% -----------------------------------------------------
\begin{titlepage}

\thispagestyle{empty}

\begin{center}

\vspace*{1cm}

{\large \textsc{\universidad}}

\vspace{0.2cm}

{\large \textsc{\facultad}}

\vspace{3cm}

{\Huge \textbf{\materia}}

\vspace{1cm}

{\LARGE \titulo}

\vspace{0.8cm}

\textit{Comprender el fundamento de cada medio de prueba es el primer paso para razonar como un jurista y no solo para memorizar un código.}

\vspace{1.2cm}

\rule{0.70\textwidth}{0.5pt}

\vspace{0.5cm}

{\large Apuntes de estudio}

\vspace{0.5cm}

\rule{0.70\textwidth}{0.5pt}

\vfill

{\large \autor}

\vspace{0.3cm}

{\large \anio}

\end{center}

\vfill

\begin{flushleft}
\textit{Este es un modesto aporte a los estudiantes de ``Derecho Procesal Civil'' no pretende sustituir el material oficial de la materia.}
\end{flushleft}

\end{titlepage}

% -----------------------------------------------------
% ÍNDICE
% -----------------------------------------------------
\tableofcontents

% -----------------------------------------------------
% INTRODUCCIÓN A LA UNIDAD
% -----------------------------------------------------
\chapter*{Introducción a la unidad}
\addcontentsline{toc}{chapter}{Introducción a la unidad}

Esta unidad aborda en profundidad tres medios de prueba fundamentales del proceso civil argentino. En primer lugar, la prueba pericial: su concepto, los sujetos involucrados (el perito designado de oficio y el consultor técnico de parte), el procedimiento completo desde el ofrecimiento hasta la valoración del dictamen, las obligaciones y derechos del perito, las causales de recusación y remoción, el anticipo de gastos y las consecuencias de su incumplimiento. En segundo lugar, se analiza la prueba indiciaria y las presunciones (legales e \textit{hominis}), explicando cómo el juez puede formar su convicción a partir de indicios acreditados en el expediente cuando no existe prueba directa. Finalmente, se estudia el reconocimiento judicial como la prueba más directa, su concepto, su procedimiento y su vinculación con las facultades del artículo 475 del Código Procesal Civil y Comercial de la Nación (CPCCN).

\textbf{Normas principales de la unidad:} arts. 459, 163 inc.\ 5, 386, 475, 476, 477 y 479 del CPCCN, y art.\ 243 del Código Penal.

\begin{note}
Puede pensarse en un juicio por accidente de tránsito: el juez no sabe si el automóvil tenía los frenos en mal estado (lo que requiere un perito mecánico), y tampoco existen testigos directos del momento de la colisión, de modo que resulta necesario acudir a indicios (marcas de frenado, velocidad, comportamiento previo del conductor). Finalmente, el juez puede trasladarse personalmente al lugar del hecho (reconocimiento judicial). Estos tres instrumentos —prueba pericial, prueba indiciaria y reconocimiento judicial— conforman el sistema que permite al juez conocer hechos técnicos o no presenciados y dictar sentencia con convicción fundada.
\end{note}

% =====================================================
% PARTE 1
% =====================================================
\chapter{La prueba pericial: concepto, sujetos y procedimiento}

\section{Introducción y finalidad de la prueba en el proceso}

Las pruebas tienen por finalidad acreditar, en el marco del proceso, la ocurrencia de un hecho de interés para la litis: un hecho controvertido, conducente, que fue fijado por el juez en la audiencia preliminar. Quien tiene la carga de acreditar ese hecho debe hacerlo a través de los medios de prueba disponibles.

Todo hecho deja un registro, una huella en las cosas o en las personas; para reproducirlo en el proceso es necesario consultar la fuente correspondiente y hacerlo a través de los medios de prueba pertinentes. Hasta este punto del programa se habían estudiado ya las pruebas documental, informativa, testimonial y confesional, todas ellas orientadas a reconstruir hechos para formar la convicción del juez. Dado que el juez no puede percibir por sí mismo la totalidad de los hechos controvertidos, debe valerse de la prueba que las partes ofrezcan y produzcan en la etapa probatoria.

\section{Cuándo es necesaria la prueba pericial}

Cuando para conocer un hecho controvertido no bastan los conocimientos comunes de la gente —ni siquiera los del juez—, sino que hace falta un conocimiento especial en alguna ciencia, arte, industria o actividad técnica especializada, las partes deben valerse de la prueba pericial.

\subsection{Características de la peritación}

\begin{definition}{Peritación}
La peritación es una actividad desarrollada por encargo judicial: no se realiza por iniciativa propia, sino que se ofrece como prueba en los escritos constitutivos del proceso (demanda, contestación, reconvención, contestación de excepciones previas), quedando a cargo del juez su ordenamiento. Es realizada por una persona ajena a las partes y al juez, con el fin de hacer conocer a este último cuestiones técnicas o científicas que escapan a su conocimiento común.
\end{definition}

\section{Quiénes son los peritos}

Si la profesión respectiva está reglamentada, son peritos quienes cuentan con título habilitante para ejercerla. Si la profesión no está reglamentada, se recurre a personas que posean conocimiento y amplia experiencia en la materia discutida en el proceso, aun sin título habilitante.

\subsection{Designación de los peritos}

Una vez ofrecida la prueba pericial y admitida por el juez, este designa y sortea un perito de la lista que llevan las cámaras de cada fuero (ingenieros, arquitectos, contadores, agrimensores, médicos de distintas especialidades, etc.), quienes se inscriben para actuar como peritos judiciales y son sorteados en el momento procesal oportuno.

\begin{note}
El término \textbf{desinsaculado} proviene de la práctica antigua de sortear al perito extrayendo un papel o una bolilla de un saco (\textit{insacular}); de allí la expresión, aún vigente, de que el perito ``ha sido desinsaculado''.
\end{note}

El sorteo se explica porque, por regla general, interviene un perito único designado de oficio por el juez cuando existen razones que justifican la intervención de una persona con conocimientos especiales. El perito es entonces un tercero ajeno al proceso, idóneo y, en su caso, poseedor del título habilitante necesario para esclarecer cuestiones técnicas, científicas, artísticas o industriales. Su función es la de un auxiliar imparcial del juez, sin más interés que el esclarecimiento de los hechos para que el juez pueda aplicar la norma correspondiente.

\section{El consultor técnico (perito de parte)}

Así como el letrado patrocinante asiste a la parte en el aspecto jurídico, la parte puede contar con un asistente en materias científicas o técnicas: el consultor técnico o perito de parte, que tanto el actor como el demandado tienen derecho a designar, incluso en los escritos constitutivos del proceso.

\begin{definition}{Consultor técnico}
El consultor técnico es un perito de parte que puede y debe ser parcial, en tanto es contratado por su cliente para fundamentar científica o técnicamente la cuestión en el sentido que lo favorezca. Se diferencia así del perito de oficio, que debe ser absolutamente imparcial y objetivo.
\end{definition}

El perito puede basarse en estudios realizados en instituciones especializadas ajenas al expediente judicial e incluso valerse de la colaboración de otros profesionales, sin que ello lo exima de elaborar personalmente el dictamen.

\subsection{Diferencias entre el perito de oficio y el consultor técnico}

\begin{table}[H]
\centering
\small
\begin{tabularx}{\textwidth}{>{\raggedright\arraybackslash}p{0.22\textwidth} X X}
\toprule
\textbf{Criterio} & \textbf{Perito (de oficio)} & \textbf{Consultor técnico} \\
\midrule
Designación & Por el juez, de oficio, por sorteo & Por la parte (actor o demandado) \\
Imparcialidad & Debe ser absolutamente imparcial y objetivo & Puede y debe ser parcial; actúa en favor de su cliente \\
Obligatoriedad & Obligado a aceptar el cargo (con matices) & Facultativo; la parte puede removerlo en cualquier momento \\
Dictamen / presentación & Obligado a presentar dictamen en tiempo y forma & No está obligado a presentar nada \\
Recusación & Puede ser recusado con causa (art.\ 17 CPCCN) & No puede ser recusado (se asume su parcialidad) \\
Honorarios & A cargo de la condena en costas & A cargo de la condena en costas \\
Base del dictamen & Expediente, estudios propios y apoyaturas técnicas & Expediente y asesoría al abogado y a la parte \\
\bottomrule
\end{tabularx}
\caption{Perito de oficio y consultor técnico}
\end{table}

\section{Actividad de la contraparte al ofrecer la prueba pericial}

Quien ofrece la prueba pericial debe indicar la especialidad del perito y los puntos de pericia, pudiendo designar en el mismo acto a su consultor técnico. Ofrecida la pericia por el actor, se corre traslado al demandado junto con el traslado de la demanda (conforme al artículo que el material identifica como 468, mencionado también como 478) para que este pueda:

\begin{enumerate}[label=\alph*)]
    \item impugnar la procedencia de la prueba pericial por no requerirse conocimientos especiales conforme al art.\ 459;
    \item manifestar desinterés en la pericia;
    \item impugnar la especialidad del perito propuesto;
    \item impugnar los puntos de pericia por considerarlos errados o improcedentes;
    \item proponer otros puntos de pericia distintos;
    \item designar su propio consultor técnico.
\end{enumerate}

\begin{formula}{Art.\ 459 CPCCN --- Procedencia de la prueba pericial}
Procederá la prueba pericial cuando la apreciación de los hechos controvertidos requiera conocimientos especiales en alguna ciencia, arte, industria o actividad técnica especializada.
\end{formula}

Si, pese a la impugnación de la contraria, el juez hace lugar a la pericia pero esta no resulta un elemento determinante en la sentencia, quien impugnó o manifestó desinterés no queda obligado al pago de las costas de la pericia. Si la contraparte impugna los puntos de pericia y ofrece otros distintos, corresponde correr traslado —por nota, no por cédula— a quien ofreció la prueba. Esta dinámica se aplica en forma simétrica según quien ofrezca la pericia, con la diferencia de que el traslado al demandado de la pericial ofrecida por el actor se notifica por cédula, mientras que el traslado al actor de la ofrecida por el demandado corre por nota.

\section{La audiencia preliminar y la designación del perito}

En la audiencia preliminar el juez debe resolver todo planteo relativo a la especialidad del perito, las impugnaciones sobre los puntos de pericia, y fijar la especialidad si las partes no se pusieron de acuerdo. Si las partes coinciden, pueden designar ellas mismas al perito, sin perjuicio de que el juez pueda cuestionar su idoneidad si carece de título habilitante. El juez conserva también la facultad de agregar o modificar los puntos de pericia propuestos.

De no existir acuerdo entre las partes, el juez determina la especialidad, fija los puntos de pericia y designa de oficio al perito, por sorteo, de entre los integrantes de la lista de la cámara del fuero correspondiente.

\section{Recusación del perito}

A partir de la designación en la audiencia preliminar corre un plazo de cinco días para que las partes recusen al perito, si lo estiman necesario.

\begin{formula}{Art.\ 17 CPCCN --- Causales de recusación con causa (aplicables al perito)}
Son causales de recusación del perito las mismas previstas para la recusación con causa del juez: tener manifiesto interés en el pleito, ser amigo íntimo o enemigo de alguna de las partes, haber sido denunciado por alguna de ellas, entre otras causales equivalentes.
\end{formula}

A estas causales se suman otras específicas de la cuestión pericial: la falta de título habilitante, cuando la profesión está reglamentada, y la incompetencia del perito, cuando su título no corresponde a la especialidad requerida (por ejemplo, la designación de un médico clínico cuando se necesita un traumatólogo o un psiquiatra, o de un ingeniero cuando corresponde un agrimensor).

Si el perito acepta las causales invocadas, el juez lo remueve directamente y designa a otro de oficio. Si las rechaza, se forma un incidente que resuelve el juez, cuya resolución es irrecurrible en cuanto a la procedencia de la recusación. Cabe destacar que la notificación de la designación del perito se practica por nota, por lo que las partes deben estar atentas para no dejar vencer el plazo de recusación.

\section{Aceptación del cargo, obligaciones y sanciones del perito}

El perito debe presentarse y aceptar efectivamente el cargo ante el Oficial Primero para intervenir válidamente en el proceso. La aceptación es facultativa —no constituye una carga pública—, aunque el rechazo reiterado e injustificado del cargo puede acarrear la exclusión del perito de la lista de la cámara correspondiente.

Aceptado el cargo, el perito asume el deber de comparecer cada vez que sea citado, presentar el dictamen en tiempo y forma, y concurrir a las audiencias para dar las explicaciones que correspondan.

\begin{formula}{Art.\ 243 del Código Penal --- Sanción penal al perito}
El perito legalmente citado que no compareciere o rehusare dar las explicaciones que le fueren requeridas será penado con prisión de quince días a un mes e inhabilitación especial de un mes a un año.
\end{formula}

Además de la sanción penal, el perito puede ser responsabilizado por los daños y perjuicios que ocasione su incumplimiento, reclamables por vía incidental.

\section{Anticipo para gastos}

Dentro de los tres días de aceptado el cargo, el perito puede solicitar un anticipo para gastos cuando la pericia requiera materiales, insumos o traslados de cierta entidad. Este anticipo no constituye un aumento encubierto de honorarios: el perito tiene el deber de rendir cuenta de los gastos realizados y de su destino en el dictamen correspondiente, de modo que las partes puedan controlarlos.

\begin{formula}{Régimen del anticipo para gastos}
La resolución que fija el anticipo para gastos solo es susceptible de recurso de reposición —no de apelación ante la cámara—. Firme la resolución, la parte obligada debe depositar el anticipo dentro de los cinco días de notificada por cédula. Si no lo deposita oportunamente, se la tiene por desistida de la prueba pericial y pierde la prueba.
\end{formula}

\section{Remoción del perito}

Son causales de remoción del perito: su renuncia durante el proceso; que, tras aceptar el cargo, rehúse presentarse a dar el dictamen; o que no lo presente oportunamente, aun habiendo solicitado ampliación del plazo, cuando la demora no esté justificada. En caso de remoción, el perito debe afrontar, si corresponde, los daños ocasionados, y es pasible de las sanciones civiles y penales pertinentes.

El plazo para presentar el dictamen, salvo disposición en contrario, es de quince días. Corresponde a la parte que ofreció la pericia realizar las diligencias necesarias para evitar la caducidad de la prueba, siendo fundamental notificar al perito de su designación para que se presente y acepte el cargo.

% =====================================================
% PARTE 2
% =====================================================
\chapter{Práctica de la pericia: pasos, dictamen y valoración}

\section{Presencia de las partes en la realización de la pericia}

Si las partes desean presenciar la realización de la prueba pericial, deben hacérselo saber al perito en el expediente.

\begin{important}
El perito tiene el deber de informar, con antelación suficiente, el lugar, día y hora en que llevará a cabo la pericia, de modo que las partes —en especial los consultores técnicos— puedan concurrir y formular las observaciones que estimen pertinentes. Si las partes solicitaron presenciar la pericia y el perito omite informar el día de su realización, la pericia resulta nula por violación del principio de bilateralidad de la audiencia, de raigambre constitucional.
\end{important}

\section{Los pasos para realizar la prueba pericial}

\subsection{Paso 1: Determinación del tema pericial}

El perito debe tener en claro qué cuestiones debe resolver, según surjan de los puntos de pericia. Puede requerir explicaciones al juez y a las partes al respecto, e incluso advertir que la cuestión excede su especialidad, en cuyo caso debe indicar cuál sería el perito de la especialidad correcta.

\subsection{Paso 2: Observación del campo de trabajo}

El perito determina cuál es su campo de trabajo y qué elementos necesita para llevar a cabo la pericia, con facultad de concurrir a los lugares pertinentes y recurrir a entidades especializadas para satisfacer los requerimientos que surjan, recolectando los elementos que luego deberá acompañar en el dictamen.

\subsection{Paso 3: Elección de los medios}

\begin{important}
La elección de los medios que el perito emplee para llevar a cabo su tarea puede dar lugar, posteriormente, a la impugnación o a la nulidad de la pericia si compromete su entidad científica o técnica.
\end{important}

El perito puede valerse de exámenes de personas, de cosas o de estudios especializados. En pericias que requieren estudios complejos —tomografías, resonancias, estudios de biología molecular— que exceden el propio saber del experto, se produce un acercamiento entre la prueba pericial y la denominada prueba científica, que para el autor Falcón debería constituir un medio probatorio autónomo. El perito puede integrar su pericia con exámenes requeridos a otras instituciones sin que ello desnaturalice su rol, aunque la frontera entre ambos medios de prueba resulta de dudosa delimitación, en tanto el conocimiento científico avanza de manera constante.

\subsection{Paso 4: Agregación y preservación de los elementos}

El perito debe preservar los medios de los que se valió para acompañarlos al expediente. En el caso de la prueba informática, por ejemplo, debe preservar y rescatar adecuadamente la evidencia digital (mensajes, archivos) para asegurar su integridad al momento de acompañarla como fundamento de sus conclusiones.

\subsection{Paso 5: Realización de las operaciones técnicas y emisión del dictamen}

El perito realiza las operaciones técnicas necesarias para evacuar los puntos de pericia y emite su dictamen, que debe ser técnico y científico, acompañado de las pruebas necesarias.

\section{La prueba documental que acompaña el perito}

El perito no puede acompañar prueba documental supliendo la negligencia de la parte que la tenía a su disposición. La prueba documental debe acompañarse, en principio, junto con la demanda, la contestación, la reconvención o la contestación de excepciones, salvo que se trate de documentos nuevos o recientemente conocidos, en cuyo caso pueden acompañarse hasta el llamamiento de autos para sentencia en primera instancia (y, en segunda instancia, fundándose libremente en la expresión de agravios). Esta regla rige cuando el perito, en el desarrollo de su pericia, solicita documentación a bancos o entidades públicas que luego acompaña como fundamento de su dictamen.

\section{Notificación del dictamen y actividad de las partes}

Presentado el dictamen por escrito y con copias para las partes, estas son notificadas por cédula. A partir de allí disponen de un plazo de cinco días para:

\begin{enumerate}[label=\alph*)]
    \item impugnar la pericia por contener falsedades;
    \item plantear la nulidad de la prueba pericial (no contemplada expresamente en el código, pero disponible por vía incidental);
    \item pedir una ampliación de la pericia, si el perito no contestó todos los puntos solicitados;
    \item pedir una nueva pericia, cuando la anterior sea inconsistente y no admita corrección por vía de ampliación;
    \item pedir explicaciones al perito.
\end{enumerate}

La nulidad de la pericia puede fundarse, por ejemplo, en que se ejerció violencia sobre el perito, en que este perdió sus facultades mentales, o en que el dictamen presenta groseros errores de fundamentación que no admiten corrección mediante un pedido de explicaciones. Si prospera la nulidad, debe realizarse una nueva pericia por otro experto, no por el perito que ya intervino.

\begin{note}
Si la impugnación por inconsistencias o falsedades no se plantea dentro de los cinco días de notificado el dictamen, puede formularse hasta el momento de los alegatos. Si tampoco se planteó en esa oportunidad, y el juez utiliza igualmente la pericia como elemento de convicción en la sentencia, la parte puede cuestionarla en la expresión de agravios al apelar la sentencia definitiva.
\end{note}

\section{El artículo 475 CPCCN: medidas complementarias}

\begin{formula}{Art.\ 475 CPCCN --- Medidas complementarias de la prueba pericial}
Inc.\ 1: planos, relevamientos, tomas fotográficas y filmaciones cinematográficas de objetos, documentos o lugares, con empleo de medios o instrumentos técnicos.\\
Inc.\ 2: exámenes científicos necesarios para el mejor esclarecimiento de los hechos controvertidos.\\
Inc.\ 3: reconstrucción de hechos, para comprobar si se produjeron o pudieron producirse de determinada manera.
\end{formula}

Las medidas del inciso 1 (planos, fotografías, películas) constituyen, en rigor, prueba documental, en tanto la prueba documental consiste en la representación del pensamiento humano en un objeto. La reconstrucción de hechos tampoco integra la prueba pericial en sentido estricto: es un medio de prueba distinto. Estas facultades del art.\ 475 se vinculan estrechamente con el reconocimiento judicial del art.\ 479.

\section{Valoración del dictamen pericial (artículo 477 CPCCN)}

\begin{formula}{Art.\ 477 CPCCN --- Valoración del dictamen pericial}
El juez valorará el dictamen pericial de acuerdo con las reglas de la sana crítica. Podrá apartarse de las conclusiones del perito cuando existan otras pruebas o razones que justifiquen tal apartamiento.
\end{formula}

El dictamen pericial no es vinculante para el juez, salvo que la ley disponga lo contrario. Para valorarlo, el juez debe ponderar cuán completo es el dictamen, cuán competentes e idóneos son los expertos que lo elaboraron, cuán coherente y lógico resulta el dictamen en sí mismo, en qué medida está fundado en cuestiones teóricas y científicas, y cuán completo es en cuanto a los elementos acompañados.

\begin{important}
En el juicio de demencia, la ley establece que el presunto insano debe ser examinado por tres médicos psiquiatras; si estos lo encuentran en sano juicio, el juez no puede declararlo demente. En este supuesto, el dictamen es absolutamente vinculante. No ocurre lo mismo a la inversa: aunque los médicos lo consideren demente, el juez —fundado en otras razones del expediente y conforme a la sana crítica— puede apartarse de esa conclusión.
\end{important}

\section{El perito árbitro: cuando el dictamen sí es vinculante}

En el proceso arbitral, lo que deciden los peritos árbitros en el ámbito de su competencia respecto de cuestiones de hecho resulta vinculante para el juez. El artículo 516 del CPCCN permite, en la etapa de ejecución de sentencia y ante cuestiones de difícil dilucidación (por ejemplo, liquidaciones de sociedades o de la sociedad conyugal), remitir la cuestión a peritos árbitros como método alternativo de resolución de conflictos, utilizando la misma regla que rige para los amigables componedores. El perito árbitro vincula al juez en cuanto a la fijación del hecho a través del laudo emitido, pero no lo limita en cuanto a la aplicación del derecho.

% =====================================================
% PARTE 3
% =====================================================
\chapter{Prueba indiciaria y reconocimiento judicial}

\section{La prueba indiciaria: concepto y necesidad}

\begin{formula}{Art.\ 476 CPCCN --- Opinión de instituciones}
Permite solicitar opiniones a universidades o a instituciones públicas o privadas respecto de cuestiones técnicas o científicas complejas. Esta previsión se engloba en lo que, para el autor Falcón, debería considerarse un medio de prueba autónomo: la prueba científica.
\end{formula}

La prueba indiciaria resulta necesaria precisamente cuando se carece de demostración acabada, a través de los demás medios probatorios (documental, informativa, pericial, reconocimiento judicial), respecto de la certeza sobre la existencia de un hecho relevante para la litis. Si los hechos están acreditados categóricamente por otros medios, no es necesario recurrir a la prueba indiciaria.

\section{Presunciones: concepto y tipos}

Existe cierta confusión terminológica entre los conceptos de indicio y presunción. Para algunos autores, la prueba indiciaria no constituye un medio de prueba autónomo sino una forma de razonamiento jurídico; para otros, en cambio, aun considerada indirectamente y sobre la base de principios lógicos, atiende a la demostración de un hecho controvertido y por ello debe ser considerada un medio de prueba.

\begin{table}[H]
\centering
\begin{tabularx}{\textwidth}{>{\raggedright\arraybackslash}p{0.42\textwidth} >{\raggedright\arraybackslash}p{0.28\textwidth} X}
\toprule
\textbf{Tipo de presunción} & \textbf{Quién la establece} & \textbf{¿Admite prueba en contrario?} \\
\midrule
Legal --- \textit{iuris et de iure} (pleno derecho) & El legislador (la ley) & No; es absoluta \\
Legal --- \textit{iuris tantum} (relativa) & El legislador (la ley) & Sí; puede desvirtuarse \\
Judicial --- \textit{hominis} & El juez & Depende de la valoración del juez \\
\bottomrule
\end{tabularx}
\caption{Tipos de presunción}
\end{table}

En la presunción legal, el legislador le atribuye peso a determinados hechos acreditados, disponiendo que el juez los tenga por ciertos y extraiga de ellos una consecuencia. En la presunción \textit{hominis}, en cambio, es el propio juez quien atribuye ese peso a los indicios.

\section{Las presunciones judiciales (\textit{hominis}): requisitos del artículo 163 inciso 5}

Las presunciones judiciales no están reguladas como medio de prueba autónomo en el capítulo correspondiente del código, sino en el artículo 163, referido a la sentencia.

\begin{formula}{Art.\ 163 inciso 5 CPCCN --- Presunciones judiciales}
Las presunciones no establecidas por la ley constituirán prueba cuando se funden en hechos reales y probados y cuando, por su número, precisión, gravedad y concordancia, produzcan convicción según la naturaleza del juicio, de conformidad con las reglas de la sana crítica.
\end{formula}

\begin{definition}{Requisitos de los indicios}
\begin{itemize}
    \item \textbf{Hechos reales y probados}: los indicios deben estar acreditados en el expediente; no pueden ser una conjetura, sino un hecho concreto y probado que, sumado a otros, contribuya a tener por acreditado el hecho presumido.
    \item \textbf{Numerosos}: debe existir pluralidad de indicios.
    \item \textbf{Precisos}: deben admitir una única interpretación razonable.
    \item \textbf{Graves}: deben tener peso propio para orientar hacia el hecho de que se trate.
    \item \textbf{Concordantes}: deben ser convergentes entre sí, formando una cadena sin lagunas ni contradicciones.
\end{itemize}
\end{definition}

\section{La conducta procesal de las partes como indicio}

El artículo 163 también dispone que la conducta de las partes puede constituir un indicio: una conducta omisiva, la destrucción de documentación, la falta de contestación de la demanda o la actuación de mala fe a lo largo del proceso pueden, sumadas a otros elementos probatorios, contribuir a la convicción del juez.

\begin{important}
El juez no podría fundar una condena únicamente en la conducta procesal de una de las partes; sin embargo, reunidos otros elementos probatorios y verificada una conducta desleal, omisiva o inequitativa, dicha conducta puede constituir un indicio adicional para fallar en contra de quien la asume.
\end{important}

\section{Las reglas de la sana crítica (artículo 386 CPCCN)}

\begin{formula}{Art.\ 386 CPCCN --- Regla de valoración de la prueba}
Salvo disposición legal en contrario, los jueces formarán su convicción respecto de la prueba de conformidad con las reglas de la sana crítica. No tendrán el deber de expresar en la sentencia la valoración de todas las pruebas producidas, sino únicamente de las que fueren esenciales y decisivas para el fallo de la causa.
\end{formula}

Las reglas de la sana crítica constituyen un razonamiento que el juez debe seguir, libre de error y de vicio, para determinar cuál de las posiciones asumidas por las partes en el proceso resulta correcta. Se trata de un método científico que el autor Falcón sistematizó en ocho reglas, vinculado con las reglas biológicas de la experiencia y de la razón, sistematizadas de modo que el juez pueda demostrar, de forma concreta, lógica y razonada, cómo formó su convencimiento, permitiendo que ese razonamiento sea controlado por las partes a los fines de una eventual impugnación.

\begin{note}
Cuando el artículo 386 dispone ``salvo disposición legal en contrario'', remite al principio de que, si el legislador le puso un ``precio'' a determinada prueba (prueba legal tasada), el juez no puede apartarse de ella y debe comenzar por esas pruebas legales tasadas. Fuera de ese supuesto, el juez forma su convicción a través de las reglas de la sana crítica.
\end{note}

\section{El reconocimiento judicial (artículo 479 CPCCN)}

\begin{definition}{Reconocimiento judicial}
El reconocimiento judicial es la prueba más directa y fiable del proceso civil: el juez se apersona y percibe directamente los hechos, sin que medie para su representación ni cosas ni personas. A diferencia del testigo, que declara sobre lo percibido por sus sentidos, o del perito, que extrae conclusiones a partir de elementos técnicos, es el propio juez quien toma conocimiento directo en el lugar de los hechos.
\end{definition}

\begin{formula}{Art.\ 479 CPCCN --- Reconocimiento judicial}
El juez podrá ordenar, de oficio o a petición de parte, el reconocimiento judicial de lugares o de cosas. Se interpreta que también comprende a personas, como en los procesos de restricción de la capacidad, en los que el juez tiene entrevista con el presunto incapaz. Cuando el conocimiento exceda al propio del juez, este puede valerse de un perito para que interprete y señale las circunstancias en juego.
\end{formula}

El reconocimiento judicial es una prueba \textit{privativa} del juez: puede ser ofrecida por una parte u ordenada de oficio, pero es el juez quien decide la utilidad que le reportará en función del conocimiento que necesita para alcanzar certeza sobre determinados hechos. En el acto del reconocimiento, el juez puede hacer comparecer a quienes considere pertinentes, requerir la concurrencia de testigos y valerse también de las medidas previstas en el art.\ 475 (planos, croquis, fotografías, filmaciones, reconstrucción de hechos) para aprehender adecuadamente los hechos que deberá fijar en el proceso.

\section{Forma de la diligencia del reconocimiento judicial}

El juez debe comparecer personalmente, junto con las personas del tribunal que considere pertinentes, y debe labrarse un acta que deje constancia de las circunstancias percibidas.

\begin{important}
La narración que el juez efectúe de sus percepciones debe ser lo más objetiva posible, puesto que la etapa probatoria no es el momento de valorar la prueba, so riesgo de incurrir en prejuzgamiento. El acta debe dejar constancia objetiva de las aclaraciones y observaciones formuladas por las partes, los consultores técnicos, los peritos y, eventualmente, los testigos, así como de las dificultades presentadas durante la diligencia.
\end{important}

% =====================================================
% CUADRO CORNELL DE REPASO
% =====================================================
\chapter{Cuadro Cornell de repaso --- Unidad 15}

El siguiente cuadro sintetiza, en formato Cornell, los conceptos centrales desarrollados en la unidad sobre prueba pericial, prueba indiciaria y reconocimiento judicial.

\begin{longtable}{
    >{\raggedright\arraybackslash}p{0.30\textwidth}
    >{\raggedright\arraybackslash}p{0.60\textwidth}
}
\caption{Cuadro Cornell --- Unidad 15: Prueba Pericial, Prueba Indiciaria y Reconocimiento Judicial} \\
\toprule
\textbf{Preguntas / palabras clave} & \textbf{Notas / respuestas} \\
\midrule
\endfirsthead

\multicolumn{2}{c}{\textit{(continuación)}} \\
\toprule
\textbf{Preguntas / palabras clave} & \textbf{Notas / respuestas} \\
\midrule
\endhead

\midrule
\multicolumn{2}{r}{\textit{Continúa en la página siguiente}} \\
\endfoot

\bottomrule
\endlastfoot

\textbf{¿Cuándo procede la prueba pericial?} &
Cuando la apreciación de los hechos controvertidos requiere conocimientos especiales en alguna ciencia, arte, industria o actividad técnica especializada que exceden la cultura media del juez (art.\ 459 CPCCN). \\

\textbf{¿Qué es la peritación?} &
Es una actividad desarrollada por encargo judicial, realizada por un tercero ajeno a las partes y al juez, con conocimientos especiales, para hacer conocer al juez cuestiones técnicas o científicas que escapan a su conocimiento común. \\

\textbf{¿Cómo se designa al perito?} &
Por sorteo (desinsaculación) entre los integrantes de la lista de peritos de la cámara del fuero correspondiente. La especialidad y los puntos de pericia se fijan en la audiencia preliminar. \\

\textbf{¿Diferencia clave entre perito y consultor técnico?} &
El perito es designado de oficio, debe ser imparcial y está obligado a presentar el dictamen. El consultor técnico es de parte, puede y debe ser parcial, y su intervención es facultativa; puede ser removido por la parte que lo designó en cualquier momento. \\

\textbf{¿Cuáles son las causales de recusación del perito?} &
Las mismas del art.\ 17 CPCCN previstas para la recusación del juez (interés en el pleito, amistad íntima, enemistad, etc.), más las específicas: falta de título habilitante y falta de la especialidad requerida para los puntos de pericia. \\

\textbf{¿Qué pasa si el perito no acepta el cargo reiteradamente?} &
Puede ser removido de la lista de la cámara correspondiente, para desalentar la conducta de elegir casos en función de los honorarios esperados. \\

\textbf{¿Qué consecuencias tiene no pagar el anticipo para gastos?} &
Si la parte obligada no deposita el anticipo dentro de los cinco días de notificada por cédula, se la tiene por desistida de la prueba pericial y pierde la prueba (la resolución solo admite recurso de reposición). \\

\textbf{¿Qué es el principio de bilateralidad en la pericia?} &
Si las partes solicitaron presenciar la pericia, el perito debe informarles con antelación suficiente el lugar, día y hora. De no hacerlo, la pericia es nula por violación del principio de bilateralidad, de raigambre constitucional. \\

\textbf{¿Cuáles son los pasos del perito para realizar la pericia?} &
1) Determinación del tema pericial. 2) Observación del campo de trabajo y recolección de elementos. 3) Elección de los medios (con cautela, pues puede generar nulidad). 4) Agregación y preservación de los elementos. 5) Realización de operaciones técnicas y emisión del dictamen. \\

\textbf{¿Qué puede hacer la parte en los cinco días posteriores al dictamen?} &
Impugnar por falsedades o inconsistencias, plantear la nulidad, pedir ampliación, pedir nueva pericia o pedir explicaciones al perito. Vencido el plazo, precluyen algunas facultades (salvo la impugnación, posible hasta los alegatos). \\

\textbf{¿Es vinculante el dictamen pericial para el juez?} &
No, salvo que la ley disponga lo contrario (por ejemplo, en el juicio de demencia, si los tres médicos psiquiatras dictaminan que la persona está sana, el juez no puede declararla demente). El juez valora el dictamen según las reglas de la sana crítica y puede apartarse fundadamente. \\

\textbf{¿Cuándo sí es vinculante el dictamen?} &
Cuando interviene un perito árbitro (art.\ 516 CPCCN, en etapa de ejecución de sentencia): su fijación de los hechos es vinculante para el juez, aunque este conserva la aplicación del derecho. \\

\textbf{¿Qué son las presunciones \textit{hominis}?} &
Son presunciones judiciales: el juez infiere un hecho no acreditado directamente a partir de indicios reales, probados en el expediente, numerosos, precisos, graves y concordantes (art.\ 163 inc.\ 5 CPCCN). \\

\textbf{Diferencia: \textit{iuris et de iure} vs.\ \textit{iuris tantum}} &
\textit{Iuris et de iure} (pleno derecho): no admite prueba en contrario, la conclusión es inexorable. \textit{Iuris tantum} (relativa): admite prueba en contrario; la presunción cede si se prueba lo contrario. \\

\textbf{¿Qué es el reconocimiento judicial?} &
La prueba más directa y fiable: el juez percibe personalmente y sin intermediarios los hechos en el lugar donde ocurrieron (lugares, cosas y personas). Se labra un acta objetiva. El juez puede valerse de peritos para interpretar lo observado y de las facultades del art.\ 475 CPCCN. \\

\textbf{¿Por qué el acta del reconocimiento judicial debe ser objetiva?} &
Porque en la etapa probatoria el juez no debe valorar la prueba todavía, ya que incurriría en prejuzgamiento. El acta debe dejar constancia objetiva de lo percibido y de las aclaraciones de partes, consultores técnicos, peritos y testigos. \\

\midrule
\multicolumn{2}{p{0.92\textwidth}}{
\textbf{Resumen:} La unidad desarrolla tres pilares del sistema probatorio civil argentino. Primero, la \textbf{prueba pericial}, medio especializado al que se recurre cuando los hechos controvertidos requieren conocimientos científicos, técnicos, artísticos o industriales que exceden la cultura media del juez: interviene un perito, tercero imparcial designado de oficio por sorteo, distinto del consultor técnico (perito de parte, parcial, no obligado a comparecer). El procedimiento sigue etapas sucesivas —ofrecimiento con puntos de pericia, traslado, audiencia preliminar, designación por sorteo, posible recusación en cinco días, aceptación del cargo, anticipo de gastos, realización de la pericia con aviso previo si fue solicitado, presentación del dictamen y notificación por cédula— tras lo cual las partes disponen de cinco días para impugnar, pedir la nulidad, la ampliación, una nueva pericia o explicaciones. El juez valora el dictamen según las reglas de la sana crítica (art.\ 386 CPCCN) y solo excepcionalmente la ley lo torna vinculante (por ejemplo, en el juicio de demencia), salvo en el arbitraje, donde el perito árbitro sí vincula al juez en cuestiones de hecho. Segundo, la \textbf{prueba indiciaria} (presunciones \textit{hominis}, art.\ 163 inc.\ 5 CPCCN): a falta de prueba directa, el juez puede inferir un hecho controvertido a partir de indicios reales, probados, numerosos, precisos, graves y concordantes, lo que se distingue de las presunciones legales (\textit{iuris et de iure}, que no admiten prueba en contrario, e \textit{iuris tantum}, que sí la admiten). Tercero, el \textbf{reconocimiento judicial} (art.\ 479 CPCCN): el juez percibe directamente lugares, cosas y personas sin intermediarios, es la prueba más directa y fiable, puede valerse de peritos para interpretar lo observado y de las facultades del art.\ 475 CPCCN. Todo el andamiaje probatorio de la unidad sirve a un único fin: formar la convicción del juez para que pueda dictar una sentencia fundada y justa.
} \\

\end{longtable}

\end{document}
